% TOP_PROBLEM: {"id": 3126, "title": "Lovász Path Removal Conjecture", "category_id": 3, "category": {"id": 3, "name": "graph_theory", "display_name": "Graph Theory", "description": "Problems involving graphs, networks, and their properties.", "slug": "graph-theory", "order_index": 3, "created_at": "2026-07-31T15:26:25.671Z"}, "status": "open", "problem_number": "OPG-46629", "set_id": 12}
% FIRST_POSTED: 2026-10-01
% ENTRY_KIND: statement_only
% UnsolvedMath ID 3126; source catalogue code OPG-46629
% !TeX program = lualatex
% Definition and sources only; this file is not an LLM solution attempt.
\documentclass[11pt,a4paper]{article}
\usepackage[margin=25mm]{geometry}
\usepackage{fontspec}
\setmainfont{lmroman10-regular.otf}[BoldFont=lmroman10-bold.otf,
 ItalicFont=lmroman10-italic.otf,BoldItalicFont=lmroman10-bolditalic.otf]
\usepackage{amsmath,amssymb,mathtools}
\usepackage{unicode-math}
\setmathfont{latinmodern-math.otf}
\AtBeginDocument{\let\setminus\smallsetminus}
\usepackage{xurl,hyperref}
\hypersetup{colorlinks=true,urlcolor=blue}
\setcounter{secnumdepth}{0}
\setlength{\parindent}{0pt}
\setlength{\parskip}{5pt}
\setlength{\emergencystretch}{3em}
\begin{document}
\section*{Lovász Path Removal Conjecture}\label{3126}
{\small UnsolvedMath ID 3126 · OPG-46629 · Graph Theory · OpenGarden}
\subsection{Definitions and mathematical statement}
A finite simple undirected graph is $k$-vertex-connected, for an integer
$k\ge1$, if it has at least $k+1$ vertices and remains connected after
deletion of any set of fewer than $k$ vertices. A path
$P=(v_0,\ldots,v_t)$ is induced if its vertices are distinct and the
only edges of the graph between these vertices are the consecutive
path edges $v_{i-1}v_i$.

The conjecture asks whether there exists a function
$f:\mathbb N\to\mathbb N$ such that, for every $k\ge1$, every
$f(k)$-vertex-connected graph $G$, and every two distinct vertices
$x,y\in V(G)$, there is an induced path $P$ from $x$ to $y$ for which
\[
 G-V(P)\text{ is }k\text{-vertex-connected}.
\]
Here $G-V(P)$ is the induced subgraph on all vertices outside the path.
All vertices of $P$, including both prescribed endpoints, are deleted.
The retained graph must still have at least $k+1$ vertices under the
connectivity convention above.

The function may grow with $k$, but cannot depend on $G$, its order,
or the choice of endpoints. The conjecture asks only that some finite
threshold exists for each $k$; determining the best possible threshold
is a sharper question.

Deleting just the edges of $P$ is a different and weaker requirement,
since that operation retains the path vertices and their other incident
edges. The path must be induced in the original graph, and the conclusion
must work for every prescribed pair, not only for some convenient pair.
\subsection{Short English statement}
Does sufficiently high vertex-connectivity always allow an induced path between any prescribed pair whose vertex deletion leaves a $k$-connected graph?
\subsection{Sources}
\begin{itemize}
\item[S1] ULAM.AI and the UnsolvedMath Contributors, supplied problems.json, record 3126 (OPG-46629), OpenGarden collection. Catalogue identity and source transcription; CC BY 4.0.
\url{https://huggingface.co/datasets/ulamai/UnsolvedMath}.
\item[S2] Open Problem Garden, Lovász Path Removal Conjecture. Original problem and discussion; the mathematical formulation below is expanded from the supplied source record.
\url{http://www.openproblemgarden.org/op/lovasz_path_removal_conjecture}.
\end{itemize}
\end{document}
